%-------------------------------------------------------------------------------
\section{Conclusion}
\label{sec:conclusion}
We present the first study on public LDAP servers deployed on the Internet. The analysis uncovered more than 80k public LDAP servers, with 12k exposing personal data. Even worse, well over one thousand servers expose hashed or even cleartext passwords of users, resulting in a total of around 3.9 million exposed user credentials. 
Furthermore, we found that more than 10k servers expose possibly sensitive internal information, and more than 600 servers are linked to some CVE, suggesting they contain security vulnerabilities. We also found that although LDAP supports TLS, either implicit via port 686 or using StartTLS over port 389, the vast majority of LDAP servers supporting TLS have issues in the TLS configuration. This is either because it uses deprecated TLS versions or cipher suites, or because the certificate is self-signed, does not validate against common trust chains, or has expired. Interestingly, the TLS configuration of servers leaking personal data or user credentials is much worse than the TLS configuration of LDAP servers that may be used for user authentication. This indicates that servers leaking personal data or user credentials are not meant to be public.

Overall, the analysis strongly suggests that a surprisingly large number of organizations run LDAP servers leaking highly critical information, running non-optimal configuration values, and subpar TLS configurations. These LDAP servers are a treasure trove for cyberattackers, giving them important insights and possibly even access credentials to privileged users. Affected organizations need to quickly mitigate these risks by securing their LDAP server or removing them from the public Internet. We believe that this analysis sets the ground for future analyses in the security of the LDAP landscape on the Internet.

%\todo{Sebastian: summarize methodogy and discussion.}

%\paragraph{Limitations.} % Unfortunately, we are unable to paint a full picture. For instance, some LDAP servers may be accessible on nonstandard internet ports. However, scanning for this server would exceed our scanning capacities and time frame. Additionally, 
%To achieve a manageable scope, we limited the analysis to the most salient and security-relevant \acrshort{LDAP} controls and features. We acknowledge that a significant number of other \acrshort{LDAP} properties remain that also deserve scrutiny. We also limit ourselves to what can be tested automatically; manual analyses and in-depth assessments of \acrshort{LDAP} server and client implementations are out of scope. We note that \acrshort{LDAP} can be used over both TCP and UDP. We use only TCP for our scans. Furthermore, we do not consider other attack vectors on \acrshort{LDAP}, such as exploiting its UDP instances in DDoS attacks.\\

%The identification of personal data is limited by the missing context for our assessment. The identified servers are ultimately candidates that unintentionally expose personal data. As an example of the missing context, servers answering with the last name and email attributes might act intentionally as a public address book, whereby servers exposing social security numbers are probably not intentionally public. 
